\section{How a task runs}

A task's \texttt{runner} says who acts:

\begin{itemize}
  \item \texttt{tool}: no model. One tool is called with the task's input,
    and its output is the task's, as in the starter's
    \code{nanoagent.yaml}{say-hello}.
  \item \texttt{agent}: an LLM loop with the task's tools, ending with
    \texttt{submit}, which takes the task's output class, as in arXiv
    Source's \arxivcode{nanoagent.yaml}{find-and-save-paper}.
  \item \texttt{decide}: the model only chooses: the next tool or
    \texttt{submit}, an Enum or \texttt{bool} member, or an existing value
    of exactly a field's type from the task's input or an earlier tool's
    output. NanoAgent calls the tools and builds every value. Every answer
    is checked against its options, and a task with a field that cannot be
    filled this way is refused when the plugin loads.
\end{itemize}

A task's \texttt{reference}, when set, is a tool outside the task's tools.
NanoAgent runs it first, out of the model's sight, and pins its output as the
ground truth for \texttt{verify\_task}, as arXiv Source's
\texttt{daily-digest} pins \texttt{list\_daily\_papers} and
\arxivcode{verifier}{its verifier} checks the digest against it. Otherwise
it is \texttt{null}.
